\documentclass[12pt, a4paper]{article}

\usepackage[T1]{fontenc}
\usepackage[utf8]{inputenc}
\usepackage[french]{babel}
\usepackage[mathscr]{euscript}
\usepackage{lmodern}
\let\euscr\mathscr \let\mathscr\relax
\usepackage[scr]{rsfso}
\usepackage{amssymb,amsmath,amsthm,amsfonts,mathtools}
\usepackage[shortlabels]{enumitem}
\usepackage{multicol,multirow}
\usepackage{lipsum}
\usepackage{balance}
\usepackage{calc}
\usepackage[colorlinks=true,citecolor=blue,linkcolor=blue]{hyperref}
\usepackage{import}
\usepackage{xifthen}
\usepackage{pdfpages}
\usepackage{transparent}
\usepackage{tabularx}
\usepackage{listings}
\usepackage[dvipsnames]{xcolor}
\usepackage{bbm}

\newcommand{\incfig}[2][1.0]{
    \def\svgwidth{#1\columnwidth}
    \import{./figures/}{#2.pdf_tex}
}
\newcommand{\incimg}[2][1.0]{
  \includegraphics[width=#1\columnwidth]{./figures/#2}
}

\newcommand{\xor}{\oplus}
\newcommand{\floor}[1]{\left\lfloor #1 \right\rfloor}
\newcommand{\ceil}[1]{\left\lceil #1 \right\rceil}

\newlist{enumproof}{enumerate}{4}
\setlist[enumproof,1]{label=\arabic*., parsep=1em}
\setlist[enumproof,2]{label=\arabic{enumproofi}.\arabic*., parsep=1em}
\setlist[enumproof,3]{label=\arabic{enumproofi}.\arabic{enumproofii}.\arabic*., parsep=1em}
\setlist[enumproof,4]{label=\arabic{enumproofi}.\arabic{enumproofii}.\arabic{enumproofiii}.\arabic*., parsep=1em}

\lstdefinestyle{mystyle}{
  language=Python, % Set the language to Python
  commentstyle=\color{codegray}\itshape, % Python comments (#)
  keywordstyle=\color{blue},  % Python reserved keywords
  stringstyle=\color{mauve},             % Strings
  basicstyle={\ttfamily\footnotesize},   % Monospace font
  breakatwhitespace=true,         
  breaklines=true,                 
  captionpos=b,                    
  keepspaces=true,                 
  numbers=none,                    
  tabsize=2,
  showstringspaces=false,
  % Extra keyword groups for Python
  morekeywords=[2]{self, cls, __init__, __str__, __repr__}, 
  keywordstyle=[2]\color{identblue},     % Special identifiers / methods
  morekeywords=[3]{int, float, bool, str, list, dict, set, tuple, object, None, True, False},
  keywordstyle=[3]\color{codegreen},     % Built-in types
  morekeywords=[4]{print, input, open, range, len, type, isinstance, super, enumerate, zip, map, filter, sum, min, max, sorted},
  keywordstyle=[4]\color{cyan},          % Built-in functions
}
\lstset{style=mystyle}

\usepackage{ifthen}
\usepackage[landscape]{geometry}
\usepackage[shortlabels]{enumitem}

\ifthenelse{\lengthtest { \paperwidth = 11in}}
    { \geometry{top=.5in,left=.5in,right=.5in,bottom=.5in} }
	{\ifthenelse{ \lengthtest{ \paperwidth = 297mm}}
		{\geometry{top=1cm,left=1cm,right=1cm,bottom=1cm} }
		{\geometry{top=1cm,left=1cm,right=1cm,bottom=1cm} }
	}

\pagestyle{empty}
\makeatletter
\renewcommand\thesection{\arabic{section}.}
\renewcommand{\section}{\@startsection{section}{1}{0mm}%
                                {-1ex plus -.5ex minus -.2ex}%
                                {0.05ex}%x
                                {\normalfont\normalsize\bfseries}}
\renewcommand{\subsection}{\@startsection{subsection}{2}{0mm}%
                                {-1ex plus -.5ex minus -.2ex}%
                                {0.05ex}%
                                {\normalfont\small\bfseries}}
\renewcommand{\subsubsection}{\@startsection{subsubsection}{3}{0mm}%
                                {-1ex plus -.5ex minus -.2ex}%
                                {0.05ex}%
                                {\normalfont\footnotesize\bfseries}}
\newcommand{\colbreak}{\vfill\null\columnbreak}
\makeatother
\setcounter{secnumdepth}{1}
\setlength{\parindent}{0pt}
\setlength{\parskip}{0.7em}

\setlist[itemize]{itemsep=0.6ex, topsep=2pt, partopsep=0pt, parsep=0pt, leftmargin=*, label=\textbullet}
\setlist[enumerate]{itemsep=0.6ex, topsep=-2pt, partopsep=0pt, parsep=0pt, label=\roman*.}

% Things Lie
\newcommand{\kb}{\mathfrak b}
\newcommand{\kg}{\mathfrak g}
\newcommand{\kh}{\mathfrak h}
\newcommand{\kn}{\mathfrak n}
\newcommand{\ku}{\mathfrak u}
\newcommand{\kz}{\mathfrak z}
\DeclareMathOperator{\Ext}{Ext} % Ext functor
\DeclareMathOperator{\Tor}{Tor} % Tor functor
\newcommand{\gl}{\opname{\mathfrak{gl}}} % frak gl group
\renewcommand{\sl}{\opname{\mathfrak{sl}}} % frak sl group chktex 6

% More script letters etc.
\newcommand{\SA}{\mathcal A}
\newcommand{\SB}{\mathcal B}
\newcommand{\SC}{\mathcal C}
\newcommand{\SF}{\mathcal F}
\newcommand{\SG}{\mathcal G}
\newcommand{\SH}{\mathcal H}
\newcommand{\OO}{\mathcal O}

\newcommand{\SCA}{\mathscr A}
\newcommand{\SCB}{\mathscr B}
\newcommand{\SCC}{\mathscr C}
\newcommand{\SCD}{\mathscr D}
\newcommand{\SCE}{\mathscr E}
\newcommand{\SCF}{\mathscr F}
\newcommand{\SCG}{\mathscr G}
\newcommand{\SCH}{\mathscr H}

% Mathfrak primes
\newcommand{\km}{\mathfrak m}
\newcommand{\kp}{\mathfrak p}
\newcommand{\kq}{\mathfrak q}

% number sets
\newcommand{\RR}[1][]{\ensuremath{\ifstrempty{#1}{\mathbb{R}}{\mathbb{R}^{#1}}}}
\newcommand{\NN}[1][]{\ensuremath{\ifstrempty{#1}{\mathbb{N}}{\mathbb{N}^{#1}}}}
\newcommand{\ZZ}[1][]{\ensuremath{\ifstrempty{#1}{\mathbb{Z}}{\mathbb{Z}^{#1}}}}
\newcommand{\QQ}[1][]{\ensuremath{\ifstrempty{#1}{\mathbb{Q}}{\mathbb{Q}^{#1}}}}
\newcommand{\CC}[1][]{\ensuremath{\ifstrempty{#1}{\mathbb{C}}{\mathbb{C}^{#1}}}}
\newcommand{\PP}[1][]{\ensuremath{\ifstrempty{#1}{\mathbb{P}}{\mathbb{P}^{#1}}}}
\newcommand{\HH}[1][]{\ensuremath{\ifstrempty{#1}{\mathbb{H}}{\mathbb{H}^{#1}}}}
\newcommand{\FF}[1][]{\ensuremath{\ifstrempty{#1}{\mathbb{F}}{\mathbb{F}^{#1}}}}
% expected value
\newcommand{\EE}{\ensuremath{\mathbb{E}}}
\newcommand{\charin}{\text{ char }}
\DeclareMathOperator{\sign}{sign}
\DeclareMathOperator{\Aut}{Aut}
\DeclareMathOperator{\Inn}{Inn}
\DeclareMathOperator{\Syl}{Syl}
\DeclareMathOperator{\Gal}{Gal}
\DeclareMathOperator{\GL}{GL} % General linear group
\DeclareMathOperator{\SL}{SL} % Special linear group

%---------------------------------------
% BlackBoard Math Fonts :-
%---------------------------------------

%Captital Letters
\newcommand{\bbA}{\mathbb{A}}	\newcommand{\bbB}{\mathbb{B}}
\newcommand{\bbC}{\mathbb{C}}	\newcommand{\bbD}{\mathbb{D}}
\newcommand{\bbE}{\mathbb{E}}	\newcommand{\bbF}{\mathbb{F}}
\newcommand{\bbG}{\mathbb{G}}	\newcommand{\bbH}{\mathbb{H}}
\newcommand{\bbI}{\mathbb{I}}	\newcommand{\bbJ}{\mathbb{J}}
\newcommand{\bbK}{\mathbb{K}}	\newcommand{\bbL}{\mathbb{L}}
\newcommand{\bbM}{\mathbb{M}}	\newcommand{\bbN}{\mathbb{N}}
\newcommand{\bbO}{\mathbb{O}}	\newcommand{\bbP}{\mathbb{P}}
\newcommand{\bbQ}{\mathbb{Q}}	\newcommand{\bbR}{\mathbb{R}}
\newcommand{\bbS}{\mathbb{S}}	\newcommand{\bbT}{\mathbb{T}}
\newcommand{\bbU}{\mathbb{U}}	\newcommand{\bbV}{\mathbb{V}}
\newcommand{\bbW}{\mathbb{W}}	\newcommand{\bbX}{\mathbb{X}}
\newcommand{\bbY}{\mathbb{Y}}	\newcommand{\bbZ}{\mathbb{Z}}

%---------------------------------------
% MathCal Fonts :-
%---------------------------------------

%Captital Letters
\newcommand{\mcA}{\mathcal{A}}	\newcommand{\mcB}{\mathcal{B}}
\newcommand{\mcC}{\mathcal{C}}	\newcommand{\mcD}{\mathcal{D}}
\newcommand{\mcE}{\mathcal{E}}	\newcommand{\mcF}{\mathcal{F}}
\newcommand{\mcG}{\mathcal{G}}	\newcommand{\mcH}{\mathcal{H}}
\newcommand{\mcI}{\mathcal{I}}	\newcommand{\mcJ}{\mathcal{J}}
\newcommand{\mcK}{\mathcal{K}}	\newcommand{\mcL}{\mathcal{L}}
\newcommand{\mcM}{\mathcal{M}}	\newcommand{\mcN}{\mathcal{N}}
\newcommand{\mcO}{\mathcal{O}}	\newcommand{\mcP}{\mathcal{P}}
\newcommand{\mcQ}{\mathcal{Q}}	\newcommand{\mcR}{\mathcal{R}}
\newcommand{\mcS}{\mathcal{S}}	\newcommand{\mcT}{\mathcal{T}}
\newcommand{\mcU}{\mathcal{U}}	\newcommand{\mcV}{\mathcal{V}}
\newcommand{\mcW}{\mathcal{W}}	\newcommand{\mcX}{\mathcal{X}}
\newcommand{\mcY}{\mathcal{Y}}	\newcommand{\mcZ}{\mathcal{Z}}


%---------------------------------------
% Bold Math Fonts :-
%---------------------------------------

%Captital Letters
\newcommand{\bmA}{\boldsymbol{A}}	\newcommand{\bmB}{\boldsymbol{B}}
\newcommand{\bmC}{\boldsymbol{C}}	\newcommand{\bmD}{\boldsymbol{D}}
\newcommand{\bmE}{\boldsymbol{E}}	\newcommand{\bmF}{\boldsymbol{F}}
\newcommand{\bmG}{\boldsymbol{G}}	\newcommand{\bmH}{\boldsymbol{H}}
\newcommand{\bmI}{\boldsymbol{I}}	\newcommand{\bmJ}{\boldsymbol{J}}
\newcommand{\bmK}{\boldsymbol{K}}	\newcommand{\bmL}{\boldsymbol{L}}
\newcommand{\bmM}{\boldsymbol{M}}	\newcommand{\bmN}{\boldsymbol{N}}
\newcommand{\bmO}{\boldsymbol{O}}	\newcommand{\bmP}{\boldsymbol{P}}
\newcommand{\bmQ}{\boldsymbol{Q}}	\newcommand{\bmR}{\boldsymbol{R}}
\newcommand{\bmS}{\boldsymbol{S}}	\newcommand{\bmT}{\boldsymbol{T}}
\newcommand{\bmU}{\boldsymbol{U}}	\newcommand{\bmV}{\boldsymbol{V}}
\newcommand{\bmW}{\boldsymbol{W}}	\newcommand{\bmX}{\boldsymbol{X}}
\newcommand{\bmY}{\boldsymbol{Y}}	\newcommand{\bmZ}{\boldsymbol{Z}}
%Small Letters
\newcommand{\bma}{\boldsymbol{a}}	\newcommand{\bmb}{\boldsymbol{b}}
\newcommand{\bmc}{\boldsymbol{c}}	\newcommand{\bmd}{\boldsymbol{d}}
\newcommand{\bme}{\boldsymbol{e}}	\newcommand{\bmf}{\boldsymbol{f}}
\newcommand{\bmg}{\boldsymbol{g}}	\newcommand{\bmh}{\boldsymbol{h}}
\newcommand{\bmi}{\boldsymbol{i}}	\newcommand{\bmj}{\boldsymbol{j}}
\newcommand{\bmk}{\boldsymbol{k}}	\newcommand{\bml}{\boldsymbol{l}}
\newcommand{\bmm}{\boldsymbol{m}}	\newcommand{\bmn}{\boldsymbol{n}}
\newcommand{\bmo}{\boldsymbol{o}}	\newcommand{\bmp}{\boldsymbol{p}}
\newcommand{\bmq}{\boldsymbol{q}}	\newcommand{\bmr}{\boldsymbol{r}}
\newcommand{\bms}{\boldsymbol{s}}	\newcommand{\bmt}{\boldsymbol{t}}
\newcommand{\bmu}{\boldsymbol{u}}	\newcommand{\bmv}{\boldsymbol{v}}
\newcommand{\bmw}{\boldsymbol{w}}	\newcommand{\bmx}{\boldsymbol{x}}
\newcommand{\bmy}{\boldsymbol{y}}	\newcommand{\bmz}{\boldsymbol{z}}

%---------------------------------------
% Scr Math Fonts :-
%---------------------------------------

\newcommand{\sA}{{\mathscr{A}}}   \newcommand{\sB}{{\mathscr{B}}}
\newcommand{\sC}{{\mathscr{C}}}   \newcommand{\sD}{{\mathscr{D}}}
\newcommand{\sE}{{\mathscr{E}}}   \newcommand{\sF}{{\mathscr{F}}}
\newcommand{\sG}{{\mathscr{G}}}   \newcommand{\sH}{{\mathscr{H}}}
\newcommand{\sI}{{\mathscr{I}}}   \newcommand{\sJ}{{\mathscr{J}}}
\newcommand{\sK}{{\mathscr{K}}}   \newcommand{\sL}{{\mathscr{L}}}
\newcommand{\sM}{{\mathscr{M}}}   \newcommand{\sN}{{\mathscr{N}}}
\newcommand{\sO}{{\mathscr{O}}}   \newcommand{\sP}{{\mathscr{P}}}
\newcommand{\sQ}{{\mathscr{Q}}}   \newcommand{\sR}{{\mathscr{R}}}
\newcommand{\sS}{{\mathscr{S}}}   \newcommand{\sT}{{\mathscr{T}}}
\newcommand{\sU}{{\mathscr{U}}}   \newcommand{\sV}{{\mathscr{V}}}
\newcommand{\sW}{{\mathscr{W}}}   \newcommand{\sX}{{\mathscr{X}}}
\newcommand{\sY}{{\mathscr{Y}}}   \newcommand{\sZ}{{\mathscr{Z}}}


%---------------------------------------
% Math Fraktur Font
%---------------------------------------

%Captital Letters
\newcommand{\mfA}{\mathfrak{A}}	\newcommand{\mfB}{\mathfrak{B}}
\newcommand{\mfC}{\mathfrak{C}}	\newcommand{\mfD}{\mathfrak{D}}
\newcommand{\mfE}{\mathfrak{E}}	\newcommand{\mfF}{\mathfrak{F}}
\newcommand{\mfG}{\mathfrak{G}}	\newcommand{\mfH}{\mathfrak{H}}
\newcommand{\mfI}{\mathfrak{I}}	\newcommand{\mfJ}{\mathfrak{J}}
\newcommand{\mfK}{\mathfrak{K}}	\newcommand{\mfL}{\mathfrak{L}}
\newcommand{\mfM}{\mathfrak{M}}	\newcommand{\mfN}{\mathfrak{N}}
\newcommand{\mfO}{\mathfrak{O}}	\newcommand{\mfP}{\mathfrak{P}}
\newcommand{\mfQ}{\mathfrak{Q}}	\newcommand{\mfR}{\mathfrak{R}}
\newcommand{\mfS}{\mathfrak{S}}	\newcommand{\mfT}{\mathfrak{T}}
\newcommand{\mfU}{\mathfrak{U}}	\newcommand{\mfV}{\mathfrak{V}}
\newcommand{\mfW}{\mathfrak{W}}	\newcommand{\mfX}{\mathfrak{X}}
\newcommand{\mfY}{\mathfrak{Y}}	\newcommand{\mfZ}{\mathfrak{Z}}
%Small Letters
\newcommand{\mfa}{\mathfrak{a}}	\newcommand{\mfb}{\mathfrak{b}}
\newcommand{\mfc}{\mathfrak{c}}	\newcommand{\mfd}{\mathfrak{d}}
\newcommand{\mfe}{\mathfrak{e}}	\newcommand{\mff}{\mathfrak{f}}
\newcommand{\mfg}{\mathfrak{g}}	\newcommand{\mfh}{\mathfrak{h}}
\newcommand{\mfi}{\mathfrak{i}}	\newcommand{\mfj}{\mathfrak{j}}
\newcommand{\mfk}{\mathfrak{k}}	\newcommand{\mfl}{\mathfrak{l}}
\newcommand{\mfm}{\mathfrak{m}}	\newcommand{\mfn}{\mathfrak{n}}
\newcommand{\mfo}{\mathfrak{o}}	\newcommand{\mfp}{\mathfrak{p}}
\newcommand{\mfq}{\mathfrak{q}}	\newcommand{\mfr}{\mathfrak{r}}
\newcommand{\mfs}{\mathfrak{s}}	\newcommand{\mft}{\mathfrak{t}}
\newcommand{\mfu}{\mathfrak{u}}	\newcommand{\mfv}{\mathfrak{v}}
\newcommand{\mfw}{\mathfrak{w}}	\newcommand{\mfx}{\mathfrak{x}}
\newcommand{\mfy}{\mathfrak{y}}	\newcommand{\mfz}{\mathfrak{z}}


\newcommand{\mytitle}{FMA\_3TC13\_Algèbre linéaire}
\newcommand{\myauthor}{github/omgeta}
\newcommand{\mydate}{AU 26-27 S1}

\newcommand{\KK}{\mathbb{K}}
\newcommand{\rg}{\operatorname{rg}}
\newcommand{\tr}{\operatorname{tr}}
\newcommand{\im}{\operatorname{im}}
\newcommand{\Com}{\operatorname{Com}}
\newcommand{\Sp}{\operatorname{Sp}}
\newcommand{\proj}{\operatorname{proj}}
\newcommand{\Vect}{\operatorname{Vect}}
\newcommand{\norm}[1]{\left\lVert #1\right\rVert}

\begin{document}
\raggedright
\footnotesize
\begin{multicols*}{3}
\setlength{\premulticols}{1pt}
\setlength{\postmulticols}{1pt}
\setlength{\multicolsep}{1pt}
\setlength{\columnsep}{2pt}

{\normalsize{\textbf{\mytitle}}}\\
{\footnotesize{\mydate\hspace{2pt}\textemdash\hspace{2pt}\myauthor}}
%%%%%%%%%%%%%%%%%%%%%%%%%%%%%%%%%%%%%%%%%%%%%%%%%%%%%%
%                      Begin                         %
%%%%%%%%%%%%%%%%%%%%%%%%%%%%%%%%%%%%%%%%%%%%%%%%%%%%%%
\section{Notions élémentaires}
\subsection{Relations et applications}
Une relation binaire entre $X$ et $Y$ est une partie de $X\times Y$. Une \textbf{application} $f:X\to Y$ associe à chaque $x\in X$ exactement un $f(x)\in Y$ 
\begin{align*}
\im f&=\{f(x):x\in X\}\subseteq Y,\\
f^{-1}(B)&=\{x\in X:f(x)\in B\},\quad B\subseteq Y
\end{align*}
\begin{enumerate}
\item Injective : $f(x)=f(x')\Rightarrow x=x'$
  \begin{itemize}
    \item Théorème: si $g\circ f$ est injective, $f$ l'est 
  \end{itemize}
\item Surjective : $\forall y\in Y,\ \exists x\in X,\ f(x)=y$
  \begin{itemize}
    \item Théorème: si $g\circ f$ est surjective, $g$ l'est 
  \end{itemize}
\item Bijective : injective et surjective\\\hspace{4.25em}$\iff f^{-1}:Y\to X$ existe
\end{enumerate}

\subsection{Relations d'ordre et d'équivalence}
Relation d'ordre : \textbf{réflexive}, \textbf{antisymétrique} et \textbf{transitive}. Relation d'équivalence $\sim$ : \textbf{réflexive}, \textbf{symétrique} et \textbf{transitive}.
\begin{align*}
[x]&=\{y\in X:x\sim y\},\\
X/{\sim}&=\{[x]:x\in X\}
\end{align*}
Les classes d'équivalence partitionnent $X$.

\section{Corps}
Un \textbf{corps} $\KK$ admet $+,$ produit, opposés et inverses multiplicatifs des éléments non nuls (ex. $\QQ,\RR,\CC$).
\begin{align*}
a\equiv b\pmod p\iff p\mid(a-b)
\end{align*}
$\ZZ/p\ZZ$ est un corps $\iff p$ premier ; $a$ inversible modulo $p$ $\iff\operatorname{pgcd}(a,p)=1$.\\

\colbreak
\section{Systèmes linéaires}
Pour $A\in\KK^{m\times n}$, $x\in\KK^n$, $b\in\KK^m$, un système s'écrit $Ax=b$ ; sa \textbf{matrice augmentée} est $[A\mid b]$.

Opérations élémentaires:
\begin{enumerate}
\item $L_i\leftarrow L_i+\lambda L_j$, $i\ne j$
\item $L_i\leftarrow\lambda L_i$, $\lambda\ne0$
\item $L_i\leftrightarrow L_j$
\end{enumerate}
Ces opérations conservent les solutions. Deux matrices reliées par celles-ci sont \textbf{équivalentes par lignes}.

\subsection{Matrice échelonnée (réduite)}
Forme \textbf{échelonnée}:
\begin{enumerate}
\item chaque première entrée non nulle d'une ligne vaut $1$ (\emph{$1$ directeur}) ;
\item lignes nulles en bas ;
\item $1$ directeurs décalés vers la droite
\end{enumerate}
Forme \textbf{réduite} : 
\begin{enumerate}
  \setcounter{enumi}{3}
  \item chaque colonne avec un $1$ directeur est nulle ailleurs
\end{enumerate}
\begin{align*}
[A\mid b]&\xrightarrow{\text{Gauss}}\text{échelonnée}\xrightarrow{}\text{réduite}
\end{align*}

Note:
\begin{enumerate}
  \item Résolution : $[0\ \cdots\ 0\mid1]\Rightarrow$ incompatible
  \item sinon chaque colonne sans pivot représente une \textbf{variable libre}
\end{enumerate}

\subsection{Systèmes homogènes}
$Ax=0$ admet toujours $x=0$. Sur un corps infini, si $\rg A<n$, il admet une infinité de solutions :
\begin{align*}
\ker A&=\{x\in\KK^n:Ax=0\},\\
\dim\ker A&=n-\rg A
\end{align*}
Pour un système à paramètres, éviter toute division par une expression susceptible de s'annuler ; étudier les valeurs exceptionnelles.

\colbreak

\section{Matrices}
Pour $A=(a_{i,j})\in\KK^{m\times n}$ et $B=(b_1\ \cdots\ b_p)\in\KK^{n\times p}$ :
\begin{align*}
(AB)_{i,k}&=\sum_{j=1}^n a_{i,j}b_{j,k},\\
AB&=(Ab_1\ \cdots\ Ab_p)
\end{align*}
\begin{enumerate}
  \item $AB \neq BA$, en général 
  \item $A,B\ne0$ peut donner $AB=0$ 
  \item $A(BC)=(AB)C$
  \item $A(B+C)=AB+AC$ 
\end{enumerate}

\subsection{Transposée}
\begin{align*}
(A^\top)_{i,j}&=a_{j,i}
\end{align*}
\begin{enumerate}
  \item $(AB)^\top=B^\top A^\top$
  \item $(A+B)^\top=A^\top+B^\top$ 
  \item $A$ symétrique $\iff A^\top=A$ 
  \item $A$ antisymétrique $\iff A^\top = -A$
\end{enumerate}

\subsection{Inversion}
$A\in GL_n(\KK)$ si $AA^{-1}=A^{-1}A=I_n$. Pour $A=\begin{psmallmatrix}a&b\\c&d\end{psmallmatrix}$ :
\begin{align*}
A^{-1}=\frac1{ad-bc}\begin{pmatrix}d&-b\\-c&a\end{pmatrix}
\end{align*}
Par \textbf{Gauss--Jordan} :
\begin{align*}
[A\mid I_n]\xrightarrow{\text{réduction}}[I_n\mid A^{-1}]
\end{align*}
\begin{enumerate}
\item $(A^{-1})^{-1}=A$
\item $(AB)^{-1}=B^{-1}A^{-1}$
\item $(A^\top)^{-1}=(A^{-1})^\top$
\item $A$ inversible $\iff\ker A=\{0\}\iff\rg A=n$
\item $A$ inversible $\iff A\sim I_n\iff\det A\ne0$
\end{enumerate}

\colbreak
\subsection{Matrices triangulaires}
Pour $A=(a_{i,j})\in\KK^{n\times n}$ :
\begin{align*}
A\text{ sup.}&\iff a_{i,j}=0\quad(i>j),\\
A\text{ inf.}&\iff a_{i,j}=0\quad(i<j)
\end{align*}
\begin{enumerate}
\item Le produit de matrices triangulaires de même type est triangulaire
\item $A$ inversible $\iff\forall i,\ a_{i,i}\ne0$ ; son inverse est triangulaire du même type
\item $\det A=\prod_{i=1}^n a_{i,i}$ 
\item $Ax=b$ : substitution de bas en haut (sup.) ou de haut en bas (inf.), si $A$ inversible
\end{enumerate}

\subsection{Puissances}
\begin{align*}
A^0&=I_n,\\
A^k&=\underbrace{A\cdots A}_{k\text{ facteurs}},\\
A^{-k}&=(A^{-1})^k\quad(A\text{ inversible}),\\
N\text{ nilpotente}&\iff\exists r\ge1,\ N^r=0
\end{align*}
\begin{enumerate}
\item $A^pA^q=A^{p+q}$ et $(A^p)^q=A^{pq}$
\item Si $AB=BA$, $(AB)^k=A^kB^k$ et formule du binôme :
\begin{align*}
(A+B)^k&=\sum_{j=0}^k\binom{k}{j}A^{k-j}B^j
\end{align*}
\item Si $A=PDP^{-1}$, $A^k=PD^kP^{-1}$ 
\item Si $A=aI_n+N$ et $N^r=0$ :
\begin{align*}
A^k&=\sum_{j=0}^{\min(k,r-1)}
\binom{k}{j}a^{k-j}N^j
\end{align*}
\end{enumerate}

\colbreak

\subsection{Trace}
Pour $A=(a_{i,j})\in\KK^{n\times n}$ :
\begin{align*}
\tr A&=\sum_{i=1}^n a_{i,i}
\end{align*}
\begin{enumerate}
\item $\tr(A+B)=\tr A+\tr B$
\item $\tr(\lambda A)=\lambda\tr A$
\item $\tr(AB)=\tr(BA)$
\item $\tr(A^\top)=\tr A$
\end{enumerate}


\subsection{Décomposition LU}
Sous les hypothèses des pivots principaux non nuls :
\begin{align*}
A&=LU,\quad L\text{ triang. inférieure},\\
&\hphantom{=LU,} U\text{ triang. supérieure}
\end{align*}
où $U$ s'obtient par Gauss sans échanges et $L$ conserve les multiplicateurs.
\begin{align*}
Ax=b\iff Ly=b,\quad Ux=y
\end{align*}
\begin{enumerate}
\item Résoudre $Ly=b$ de haut en bas
\item Résoudre $Ux=y$ de bas en haut
\end{enumerate}

\subsection{Chaînes de Markov}
$T=(p_{i,j})$ : matrice de transition ($p_{i,j}\ge0$, $\sum_i p_{i,j}=1$ pour chaque colonne)
\begin{align*}
x_{k+1}&=Tx_k, &x_k&=T^k x_0
\end{align*}
Si une distribution stationnaire existe :
\begin{align*}
Tx_\infty&=x_\infty,&\sum_i(x_\infty)_i&=1
\end{align*}

\colbreak
\section{Déterminants}
Pour $A\in\KK^{n\times n}$, le déterminant est noté $\det A$ ou $|A|$ :
\begin{align*}
\det A&=\sum_{\sigma\in\mathfrak S_n}\varepsilon(\sigma)
\prod_{i=1}^n a_{i,\sigma(i)}\\
      &= \sum_{j=1}^n a_{i,j}c_{i,j} =\sum_{i=1}^n a_{i,j}c_{i,j}\\
\det\begin{pmatrix}a&b\\c&d\end{pmatrix}&=ad-bc
\end{align*}
\begin{enumerate}
\item $\det(A^\top)=\det A$
\item $\det(AB)=\det A\det B$
\item $\det(\lambda A)=\lambda^n\det A$
\item $\det(A^{-1})=(\det A)^{-1}$ si inversible
\item $\det A=0\iff$ colonnes liées
\end{enumerate}

Effét des opérations linéaires sur le déterminant:
\begin{enumerate}
\item $L_i\leftarrow L_i+\lambda L_j$ : inchangé
\item $L_i\leftarrow\lambda L_i$ : multiplié par $\lambda$
\item $L_i\leftrightarrow L_j$ : change de signe
\end{enumerate}
Deux lignes ou colonnes égales $\Rightarrow\det A=0$

\subsection{Cofacteurs et comatrice}
\vspace{-2em}
\begin{align*}
\Com(A)&=(c_{i,j})_{i,j},\\
A\Com(A)^\top&=(\det A)I_n,\\
A^{-1}&=\frac{\Com(A)^\top}{\det A},\\
\det\Com(A)&=(\det A)^{n-1}
\end{align*}
où mineur $M_{i,j}$ efface la ligne $i$ et la colonne $j$, et  cofacteur $c_{i,j}=(-1)^{i+j}\det M_{i,j}$ 

\subsubsection{Règle de Cramer}
Si $\det A\ne0$, l'unique solution de $Ax=b$ :
\begin{align*}
x_j=\frac{\det A_j(b)}{\det A},
\end{align*}
avec $A_j(b)$ obtenu en remplaçant la colonne $j$ de $A$ par $b$.

\colbreak
\section{Espaces vectoriels}
Un espace vectoriel $V$ sur $\KK$ est stable par addition et multiplication scalaire et possède notamment un vecteur nul $0_V$ et des opposés.

\subsection{Sous-espace vectoriel}
$F\subseteq V$ est un sous-espace si :
\begin{enumerate}
\item $0_V\in F$ ;
\item $u,v\in F\Rightarrow u+v\in F$ ;
\item $u\in F,\lambda\in\KK\Rightarrow\lambda u\in F$
\end{enumerate}
Somme et somme directe de sous-espaces :
\begin{align*}
F+G&=\{f+g:f\in F,g\in G\},\\
V=F\oplus G&\iff\begin{cases}V=F+G,\\F\cap G=\{0\}\end{cases}
\end{align*}
$F\cap G$ et $F+G$ sont des sous-espaces.

\subsection{Combinaison linéaire / engendré}
Pour $u_1,\ldots,u_p\in V$ :
\begin{align*}
v&=\sum_{i=1}^p\lambda_i u_i,\\
\langle u_1,\ldots,u_p\rangle
&=\left\{\sum_{i=1}^p\lambda_i u_i:\lambda_i\in\KK\right\}
\end{align*}
Définitions d'une famille $(u_1,\ldots,u_p)$ :
\begin{align*}
\text{Génératrice}&\iff\langle u_1,\ldots,u_p\rangle=V,\\
\text{Libre}&\iff\Bigl[\sum_i\lambda_i u_i=0 \Rightarrow\forall i,\lambda_i=0\Bigr],\\
\text{Liée}&\iff\text{non libre}
\end{align*}
\begin{enumerate}
\item Si $0$ appartient à la famille, elle est liée
\item Si $p>\dim V$, elle est liée
\end{enumerate}

\colbreak
\subsection{Bases et dimension}
Une \textbf{base} est libre et génératrice ; tout vecteur s'y écrit de façon unique :
\begin{align*}
\mcB=(u_1,\ldots,u_n),\quad
v&=\sum_i\lambda_i u_i,\\
[v]_{\mcB}&=(\lambda_1,\ldots,\lambda_n)^\top
\end{align*}
\begin{enumerate}
\item $\dim V=n$ : nombre de vecteurs d'une base
\item Une famille libre se complète en une base
\item $\dim(F+G)=\dim F+\dim G-\dim(F\cap G)$\\\hfill(Grassmann)
\end{enumerate}

\subsection{Rang et noyau d'une matrice}
Pour $A=(a_1\ \cdots\ a_n)\in\KK^{m\times n}$ :
\begin{align*}
\im A&=\langle a_1,\ldots,a_n\rangle,\\
\ker A&=\{x\in\KK^n:Ax=0\},\\
\rg A&=\dim\im A
\end{align*}
\begin{enumerate}
\item $\rg A=$ nombre de pivots
\item $\rg A=\rg(A^\top)$
\item $\rg A+\dim\ker A=n$
\item $\rg(AB)\le\min(\rg A,\rg B)$
\end{enumerate}

\subsection{Changement de bases}
Pour $\mcB=(u_1,\ldots,u_n)$ et $\mcB'=(u'_1,\ldots,u'_n)$, la \textbf{matrice de passage} de $\mcB'$ vers $\mcB$ est :
\begin{align*}
P_{\mcB,\mcB'}=\bigl([u'_1]_{\mcB}\ \cdots\ [u'_n]_{\mcB}\bigr)
\end{align*}
\begin{enumerate}
\item $[v]_{\mcB}=P_{\mcB,\mcB'}[v]_{\mcB'}$
\item $P_{\mcB',\mcB}=P_{\mcB,\mcB'}^{-1}$
\item Base canonique $\mcC$ : $P_{\mcC,\mcB}=(u_1\ \cdots\ u_n)$ (coordonnées canoniques)
\end{enumerate}

\colbreak
\section{Applications linéaires}
$f:V\to W$ est \textbf{linéaire / homomorphisme} si :
\begin{align*}
f(u+v)&=f(u)+f(v),\\
f(\lambda u)&=\lambda f(u)
\end{align*}
Terminologie :
\begin{align*}
\text{Endomorphisme}&:\ V=W,\\
\text{Isomorphisme}&:\ f\text{ bijective},\\
\text{Automorphisme}&:\ V=W,\ f\text{ bijective}
\end{align*}

\subsection{Noyau, image, rang}
\begin{align*}
\ker f&=\{v\in V:f(v)=0_W\},\\
\im f&=\{f(v):v\in V\},\\
\rg f&=\dim\im f,\\
\dim V&=\dim\ker f+\rg f
\end{align*}
\begin{enumerate}
\item $f$ injective $\iff\ker f=\{0\}$
\item $f$ surjective $\iff\im f=W$
\item Si $\dim V=\dim W<\infty$, injective $\iff$ surjective $\iff$ bijective
\end{enumerate}

\subsection{Matrice d'une application}
Pour $\mcB=(u_i)$ base de $V$, $\mcD=(w_j)$ base de $W$ :
\begin{align*}
[f]_{\mcD,\mcB}
&=\bigl([f(u_1)]_{\mcD}\ \cdots\ [f(u_n)]_{\mcD}\bigr),\\
[f(v)]_{\mcD}&=[f]_{\mcD,\mcB}[v]_{\mcB},\\
[g\circ f]_{\mcE,\mcB}&=[g]_{\mcE,\mcD}[f]_{\mcD,\mcB}
\end{align*}
Pour un endomorphisme, $P=P_{\mcB,\mcB'}$ :
\begin{align*}
\boxed{[f]_{\mcB'}=P^{-1}[f]_{\mcB}P}
\end{align*}

\colbreak

\section{Espaces euclidiens}
Dans $\RR^n$, produit scalaire canonique et norme :
\begin{align*}
\langle u,v\rangle&=u^\top v=\sum_i u_iv_i,\\
\norm{v}&=\sqrt{\langle v,v\rangle},\\
\langle u,v\rangle&=\norm{u}\norm{v}\cos\theta,\\
\norm{u+v}^2&=\norm u^2+\norm v^2+2\langle u,v\rangle
\end{align*}
\begin{enumerate}
\item $|\langle u,v\rangle|\le\norm u\norm v$ \hfill(Cauchy-Schwarz)
\end{enumerate}

\subsection{Orthogonal et orthonormal}
Définitions d'orthogonalité pour $(u_i)$ :
\begin{align*}
\text{Orthogonale}&\iff\langle u_i,u_j\rangle=0\ (i\ne j),\\
\text{Orthonormée}&\iff\langle u_i,u_j\rangle=\delta_{ij},\\
F^\perp&=\{v:v\perp u,\ \forall u\in F\}
\end{align*}
\begin{enumerate}
\item Toute famille orthogonale de vecteurs est libre
\item $V=F\oplus F^\perp$ et $\dim F^\perp=\dim V-\dim F$
\end{enumerate}

\subsection{Gram--Schmidt}
D'une famille libre $(v_1,\ldots,v_p)$ :
\begin{align*}
w_1&=v_1,\\
w_k&=v_k-\sum_{i=1}^{k-1}
\frac{\langle v_k,w_i\rangle}{\norm{w_i}^2}w_i,\\
e_k&=\frac{w_k}{\norm{w_k}}
\end{align*}
$(e_1,\ldots,e_p)$ est orthonormée et engendre le même espace.

\subsection{Décomposition QR}
Pour $A$ à colonnes libres :
\begin{align*}
A&=QR, &Q^\top Q&=I,\\
R&=Q^\top A\quad\text{(triang. sup.)}
\end{align*}
Les colonnes de $Q$ viennent de Gram--Schmidt.

\colbreak
\subsection{Projection orthogonale}
Si $(e_1,\ldots,e_p)$ est une base orthonormée de $F$ :
\begin{align*}
\proj_F(v)&=\sum_{i=1}^p\langle v,e_i\rangle e_i,\\
v&=\proj_F(v)+\proj_{F^\perp}(v)
\end{align*}
Avec $Q=(e_1\ \cdots\ e_p)$ :
\begin{align*}
\proj_F(v)&=QQ^\top v,\\
P_F&=QQ^\top,\quad P_F^2=P_F=P_F^\top
\end{align*}
Avec $A$ à colonnes libres engendrant $F$ :
\begin{align*}
P_F=A(A^\top A)^{-1}A^\top
\end{align*}

\subsection{Moindres carrés}
Si $Ax=b$ est incompatible, minimiser $\norm{Ax-b}$ :
\begin{align*}
A^\top(Ax_0-b)&=0,\\
\boxed{A^\top Ax_0=A^\top b}
\end{align*}
Si les colonnes de $A$ sont libres :
\begin{align*}
x_0=(A^\top A)^{-1}A^\top b
\end{align*}
Si $A=QR$ et $R$ est inversible : $x_0=R^{-1}Q^\top b$

\subsection{Matrices orthogonales}
\vspace{-2em}
\begin{align*}
Q\in O_n(\RR)&\iff Q^\top Q=I_n
\end{align*}
\begin{enumerate}
\item $Q^{-1}=Q^\top$
\item $\det Q=\pm1$
\item $\norm{Qv}=\norm v$ et $\langle Qu,Qv\rangle=\langle u,v\rangle$
\end{enumerate}
Rotation de $\RR^2$ :
\begin{align*}
R_\theta=\begin{pmatrix}\cos\theta&-\sin\theta\\
\sin\theta&\cos\theta\end{pmatrix}
\end{align*}
Dans $\RR^3$, $u\times v$ est orthogonal à $u,v$ et $\norm{u\times v}=\norm u\norm v|\sin\theta|$.

\colbreak
\section{Diagonalisation}
\subsection{Valeurs et vecteurs propres}
Pour $A\in\KK^{n\times n}$, $v\ne0$ :
\begin{align*}
Av=\lambda v&\iff(A-\lambda I)v=0,\\
E_\lambda&=\ker(A-\lambda I),\\
\chi_A(X)&=\det(XI_n-A),\\
\Sp_\KK(A)&=\{\lambda\in\KK:\chi_A(\lambda)=0\}
\end{align*}
Multiplicités : \textbf{algébrique} = ordre de la racine $\lambda$ de $\chi_A$ ; \textbf{géométrique} = $\dim E_\lambda$
\begin{enumerate}
\item Multiplicité géométrique $\le$ multiplicité algébrique
\end{enumerate}
Avec les valeurs propres comptées avec multiplicités (dans un corps de décomposition) :
\begin{align*}
\tr A=\sum_i\lambda_i,\qquad\det A=\prod_i\lambda_i
\end{align*}

\subsection{Critère de diagonalisation}
\begin{enumerate}
\item $A$ diagonalisable $\iff$ il existe une base de vecteurs propres
\item Critère équivalent : $\displaystyle\sum_{\lambda\in\Sp_\KK(A)}\dim E_\lambda=n$
\item Des valeurs propres distinctes donnent des vecteurs propres libres
\item $\chi_A$ scindé à racines simples $\Rightarrow A$ diagonalisable
\end{enumerate}
Si $A$ est diagonalisable :
\begin{align*}
A&=PDP^{-1},\quad P=(v_1\ \cdots\ v_n),\\
D&=\operatorname{diag}(\lambda_1,\ldots,\lambda_n),\\
A^k&=PD^kP^{-1}
\end{align*}
\colbreak

\subsection{Polynômes d'endomorphismes}
Pour un polynôme $p$ :
\begin{align*}
p(X)&=\sum_{j=0}^r a_jX^j, &p(A)&=\sum_{j=0}^r a_jA^j
\end{align*}
\begin{enumerate}
\item $\chi_A(A)=0$ \hfill(Cayley-Hamilton)
\item Polynôme minimal $\mu_A$ : unitaire de plus petit degré annulant $A$
\item $\mu_A\mid\chi_A$ et divise tout polynôme annulateur
\item $A$ diagonalisable $\iff\mu_A$ scindé à racines simples
\end{enumerate}

\subsection{Réduction de Jordan}
Sur $\CC$ (ou si $\chi_A$ est scindé), $A=PJP^{-1}$ avec $J$ diagonale par blocs :
\begin{align*}
J_r(\lambda)=\begin{pmatrix}
\lambda&1&0&\cdots&0\\
0&\lambda&1&\ddots&\vdots\\
\vdots&&\ddots&\ddots&0\\
0&\cdots&0&\lambda&1\\
0&\cdots&\cdots&0&\lambda
\end{pmatrix}
\end{align*}
\begin{enumerate}
\item $A$ diagonalisable $\iff$ tous les blocs ont taille $1$
\item Un bloc de taille $r$ associé à $\lambda$ apporte le facteur $(X-\lambda)^r$ au polynôme minimal
\end{enumerate}

\subsection{Théorème spectral}
Toute matrice \textbf{symétrique réelle} $A=A^\top$ est orthogonalement diagonalisable :
\begin{align*}
A&=QDQ^\top,\quad Q^\top Q=I_n,\\
A&=\sum_{i=1}^n\lambda_i q_iq_i^\top
\end{align*}
\begin{enumerate}
\item Ses valeurs propres sont réelles
\item Des espaces propres associés à des valeurs propres distinctes sont orthogonaux
\end{enumerate}

\subsection{Valeurs singulières (SVD)}
Pour $A\in\RR^{m\times n}$ de rang $r$ :
\begin{align*}
\boxed{A=U\Sigma V^\top},\\
U^\top U=I_m,\quad V^\top V=I_n
\end{align*}
Définition des valeurs singulières non nulles (ordre décroissant) :
\begin{align*}
\sigma_i&=\sqrt{\lambda_i(A^\top A)}>0,\\
A^\top Av_i&=\sigma_i^2v_i,\\
u_i&=\frac{Av_i}{\sigma_i}\quad(i\le r),\\
A&=\sum_{i=1}^r\sigma_i u_iv_i^\top
\end{align*}
Compléter les vecteurs singuliers en bases orthonormées pour obtenir $U$ et $V$ carrées.

\colbreak
\section*{Commandes SageMath}
\subsection{Objets et calculs}
{\scriptsize\begin{tabular}{@{}l@{\hspace{2pt}}l@{}}
\texttt{QQ}, \texttt{RR}, \texttt{CC}, \texttt{SR} & corps / symboles\\
\texttt{var('a')} & variable symbolique\\
\texttt{vector(QQ,[1,2])} & vecteur exact\\
\texttt{matrix(QQ,L)} & matrice ($L$ : listes)\\
\texttt{identity\_matrix(n)} & $I_n$\\
\texttt{zero\_matrix(m,n)} & matrice nulle\\
\texttt{A.augment(b)} & matrice $[A\mid b]$\\
\texttt{e.expand()}, \texttt{e.factor()} & dév. / fact.\\
\end{tabular}}

\subsection{Matrices et systèmes}
{\scriptsize\begin{tabular}{@{}l@{\hspace{2pt}}l@{}}
\texttt{A*B}, \texttt{A\^{}k} & produit / puissance\\
\texttt{A.transpose()} & transposée\\
\texttt{A.inverse()} & inverse\\
\texttt{A.echelon\_form()} & échelonnée réduite\\
\texttt{A.solve\_right(b)} & résoudre $Ax=b$\\
\texttt{A.rank()} & rang\\
\texttt{A.determinant()} & déterminant\\
\texttt{A.trace()} & trace\\
\texttt{A.charpoly()} & polynôme car.\\
\texttt{A.minpoly()} & polynôme min.\\
\texttt{A.right\_kernel()} & noyau\\
\texttt{A.column\_space()} & image (colonnes)\\
\texttt{A.row\_space()} & espace des lignes\\
\end{tabular}}

\subsection{Réductions / espaces}
{\scriptsize\begin{tabular}{@{}l@{\hspace{2pt}}l@{}}
\texttt{A.eigenvectors\_right()} & vecteurs propres\\
\texttt{A.right\_eigenmatrix()} & $D,P$ propres\\
\texttt{A.jordan\_form()} & forme de Jordan\\
\texttt{span([u,v])} & espace engendré\\
\texttt{F.intersection(G)} & intersection\\
\texttt{F+G}, \texttt{F.dimension()} & somme / dimension\\
\texttt{u.cross\_product(v)} & produit vectoriel\\
\end{tabular}}

\textit{Pour les calculs exacts, préférer \texttt{QQ} aux flottants ; utiliser \texttt{SR} pour les paramètres symboliques.}
\end{multicols*}
%%%%%%%%%%%%%%%%%%%%%%%%%%%%%%%%%%%%%%%%%%%%%%%%%%%%%%
%                       End                          %
%%%%%%%%%%%%%%%%%%%%%%%%%%%%%%%%%%%%%%%%%%%%%%%%%%%%%%
\end{document}
